We first look at labelled transition systems (LTSs) and the modal $\mu$-calculus and what it means to verify an LTS. The definitions below are derived from \cite{Groote}.
\begin{definition}
	\label{def_lts}A labelled transition system (LTS) is a tuple $M = (S, Act, trans, s_0)$, where:
	\begin{itemize}
		\item $S$ is a set of states,
		\item $Act$ a set of actions,
		\item $trans \subseteq S \times Act \times S$ is the transition relation with $(s,a,s') \in trans$ denoted by $s \xrightarrow a s'$,
		\item $s_0 \in S$ is the initial state.
	\end{itemize}
\end{definition}
Consider the example in figure \ref{fig:coffeemachinebasiceurolts} (directly taken from \cite{FamBasedModelCheckingWithMCRL2}) of a coffee machine where we have two actions: ins (insert coin) and std (get standard sized coffee).\\
\begin{figure}[h]
	\centering
	\includegraphics[scale=0.3]{Examples/CoffeeMachine/BasicEuroLTS}
	\caption[Coffee machine LTS]{Coffee machine LTS $C$}
	\label{fig:coffeemachinebasiceurolts}
\end{figure}


\begin{definition}
	\label{def_mu_syntax}
	A modal $\mu$-calculus formula over the set of actions $Act$ and a set of variables $\mathcal{X}$ is defined by
	\[ \varphi = \top\ |\ \bot\ |\ X\ |\ \varphi \vee \varphi\ |\ \varphi \wedge \varphi\ |\ \langle a \rangle \varphi\ |\ [a]\varphi\ |\ \mu X.\varphi\ |\ \nu X.\varphi \]
	with $a \in Act$ and $X \in \mathcal{X}$. 
	
	
	No negations in the language because negations can be pushed inside to the propositions, ie. the $\top$ and $\bot$ elements.
\end{definition}
The modal $\mu$-calculus contains boolean constants $\top$ and $\bot$, propositional operators $\vee$ and $\wedge$, modal operators $\langle \, \rangle$ and $[ \, ]$ and fixpoint operators $\mu$ and $\nu$. A formula is closed when variables only occur in the scope of a fixpoint operator for that variable.

A modal $\mu$-calculus formula can be interpreted with an LTS, this results in a set of states for which the formula holds.
\begin{definition}
	\label{def_mu_sem} For LTS $(S, Act, trans, s_0)$ we inductively define the interpretation of a modal $\mu$-calculus formula $\varphi$, notation
	$\llbracket \varphi \rrbracket^\eta$, where $\eta : \mathcal{X} \rightarrow \mathcal{P}(S)$ is a logical variable valuation, as a set of states
	where $\varphi$ is valid, by:
	\begin{align*}
	&\llbracket {\top} \rrbracket^\eta &&= S\\
	&\llbracket {\bot} \rrbracket^\eta &&= \emptyset\\
	&\llbracket \varphi_1 \wedge \varphi_2 \rrbracket^\eta &&= \llbracket \varphi_1 \rrbracket^\eta \cap \llbracket \varphi_2 \rrbracket^\eta \\
	&\llbracket \varphi_1 \vee \varphi_2 \rrbracket^\eta &&= \llbracket \varphi_1 \rrbracket^\eta \cup \llbracket \varphi_2 \rrbracket^\eta\\
	&\llbracket \langle a \rangle \varphi \rrbracket^\eta &&= \{s \in S|\exists_{s' \in S} s \xrightarrow {a} s' \wedge s' \in \llbracket \varphi \rrbracket^\eta\}\\
	&\llbracket [ a ] \varphi \rrbracket^\eta &&= \{s \in S|\forall_{s' \in S} s \xrightarrow {a} s' \implies s' \in \llbracket \varphi \rrbracket^\eta\}\\
	&\llbracket \mu X. \varphi \rrbracket^\eta &&= \bigcap_{f \subseteq S}\{f | f = \llbracket \varphi \rrbracket^{\eta[X:=f]}\}\\
	&\llbracket \nu X. \varphi \rrbracket^\eta &&= \bigcup_{f \subseteq S}\{f | f = \llbracket \varphi \rrbracket^{\eta[X:=f]}\}\\
	&\llbracket X \rrbracket^\eta &&= \eta(X)
	\end{align*}
\end{definition}

Given closed formula $\varphi$, LTS $M = (S, Act, trans, s_0)$ and $s \in S$ we write $(M,s) \models \varphi$ iff $s \in \llbracket \varphi \rrbracket^\eta$ for $M$, we say that formula $\varphi$ holds for $M$ in state $s$. If formula $\varphi$ holds for $M$ in the initial state we say that formula $\varphi$ holds for $M$ and write $M \models \varphi$.

Again consider the coffee machine example (figure \ref{fig:coffeemachinebasiceurolts}) and formula $\varphi = \nu X. \mu Y. ([ins]Y \wedge [std] X)$ (taken from \cite{FamBasedModelCheckingWithMCRL2}) which states that action \textit{std} must occur infinitely often over all runs. Obviously this holds for the coffee machine, therefore we have $C \models \varphi$.

\subsection{Featured transition systems}
A \textit{featured transition system} (FTS) extends the LTS definition to express variability. It does so by introducing \textit{features} and \textit{products} into the definition. Features are options that can be enabled or disabled for the system. A product is a feature assignments, ie. a set of features that is enabled for that product. Not all products are valid, some features might be mutually exclusive and some features might always be required. To express the relation between features one can use feature diagrams as explained in \cite{Classen2013FeaturedTS}. Feature diagrams offer a nice way of expressing which feature assignments are valid, however for simplicity we will represent the collection of valid products simply with a set of feature assignments. Finally FTSs guard every transition with a boolean expression over the set of features. We have the following definition, based on \cite{Classen2013FeaturedTS}:
\begin{definition}
	\label{def_fts}A featured transition system (FTS) is a tuple $M = (S, Act, trans, s_0, N, P, \gamma)$, where:
	\begin{itemize}
		\item $S, Act, trans, s_0$ are defined as in an LTS,
		\item $N$ is a non-empty set of features,
		\item $P \subseteq \mathcal{P}(N)$ is a set of products, ie. feature assignments, that are valid,
		\item $\gamma : trans \rightarrow \mathbb{B}(N)$ is a total function, labelling each transition with a boolean expression over the features. A product $p \in \mathcal{P}(N)$ satisfying the boolean expression of transition $t$ is denoted by $p \models \gamma(t)$. The boolean expression that is satisfied by any feature assignment is denoted by $\top$, ie $p \models \top$ for any $p$.
		
		A transition $s \xrightarrow a s'$ and $\gamma((s,a,s')) = f$ is denoted by $s \xrightarrow {a | f} s'$. 
	\end{itemize}
\end{definition}
\begin{figure}[h]
\centering
\includegraphics[scale=0.3]{Examples/CoffeeMachine/FTS}
\caption[Coffee machine LTS]{Coffee machine FTS $C$}
\label{fig:coffeemachinefts}
\end{figure}

Consider the example in figure \ref{fig:coffeemachinefts} (directly taken from \cite{FamBasedModelCheckingWithMCRL2}) which shows an FTS for a coffee machine For this example we have two features $N = \{\$, \officialeuro\}$ and two valid products $P = \{\{\$\},\{\officialeuro\}\}$.

An FTS expresses the behaviour of multiple products, we can derive the behaviour of a single product by simply removing all the transitions from the FTS for which the product doesn't satisfy the feature expression guarding the transition. We call this a \textit{projection} \cite{Classen2013FeaturedTS}.

\begin{definition}
	\label{def_fts_proj}
	The projection of an FTS $M = (S, Act, trans, s_0, N, P, \gamma)$ to a product $p \in P$, noted $M_{|p}$, is the LTS $M'=(S,Act,trans', s_0)$, where $trans' = \{t \in trans\ |\ p \models \gamma(t)\}$.
\end{definition}
The coffee machine example can be projected to its two products, which results in the LTSs in figure \ref{fig:cofeemachineftsproj}.
\begin{figure}[h]
	\centering
	\begin{subfigure}{.5\textwidth}
		\centering
		\includegraphics[width=1\linewidth]{Examples/CoffeeMachine/FTSProjDollar}
		\caption[$C_{|\{\$\}}$]{$C$ projected to the dollar product: $C_{|\{\$\}}$}
		\label{fig:coffeemachineftsprojdollar}
	\end{subfigure}%
	\begin{subfigure}{.5\textwidth}
		\centering
		\includegraphics[width=1\linewidth]{Examples/CoffeeMachine/FTSProjEuro}
		\caption[$C_{|\{\$\}}$]{$C$ projected to the euro product: $C_{|\{\officialeuro\}}$}
		\label{fig:coffeemachineftsprojeuro}
	\end{subfigure}
	\caption{Projections of the coffee machine FTS}
	\label{fig:cofeemachineftsproj}
\end{figure}
\subsection{FTS verification question}
When verifiying an FTS against a model $\mu$-calculus formula $\varphi$, we are trying to answer the question: For which products in the FTS does its projection satisfy $\varphi$? Formally, given FTS $M = (S, Act, trans, s_0, N, P, \gamma)$ and modal $\mu$-calculus formula $\varphi$ we want to find $P_s \subseteq P$ such that:
\begin{itemize}
	\item for every $p \in P_s$ we have $M_{|p} \models \varphi$ and
	\item for every $p \in P \backslash P_s$ we have $M_{|p} \not\models \varphi$.
\end{itemize}
Furthermore a counterexample for every $p \in P \backslash P_s$ is preferred.